% Generated draft: requires mathematical and source review.
\documentclass[11pt]{article}
\usepackage[margin=25mm]{geometry}
\usepackage{fontspec}
\setmainfont{Latin Modern Roman}
\usepackage{amsmath,amssymb,mathtools,unicode-math}
\setmathfont{Latin Modern Math}
\usepackage{enumitem,needspace,xurl,hyperref}
\setlength{\parindent}{0pt}
\setlength{\parskip}{5pt}
\setlength{\emergencystretch}{3em}
\newcommand{\N}{\mathbb{N}}
\newcommand{\Z}{\mathbb{Z}}
\newcommand{\Q}{\mathbb{Q}}
\newcommand{\R}{\mathbb{R}}
\newcommand{\C}{\mathbb{C}}
\newcommand{\F}{\mathbb{F}}
\newcommand{\A}{\mathbb{A}}
\newcommand{\PP}{\mathbb{P}}
\newcommand{\E}{\mathbb{E}}
\newcommand{\eps}{\varepsilon}
\newcommand{\dd}{\,\mathrm{d}}
\begin{document}
\section*{UnsolvedMath ID: 700906}
\textbf{Source problem code:} catalogue entry
\subsection{Definitions and mathematical statement}
Let \(\Gamma\) denote an infinite connected graph. For a graph \(H\), write \(V(H)\) for its vertex set and \(\operatorname{Aut}(H)\) for its full automorphism group. Under the usual interpretation that primitivity refers to the action on vertices, this action is primitive if it is transitive and the only \(\operatorname{Aut}(H)\)-invariant equivalence relations on \(V(H)\) are equality and the universal relation.

The problem, proposed by V. I. Trofimov in 1992, is to characterize all such \(\Gamma\) for which there exists a sequence \((\Gamma_n)_{n\geq 1}\) of finite connected graphs satisfying
\[
\operatorname{Aut}(\Gamma_n)\text{ acts primitively on }V(\Gamma_n)
\quad\text{for every }n\geq 1,
\qquad
\Gamma_n\longrightarrow\Gamma.
\]
Here convergence must have the meaning specified in the cited paper of Trofimov, \emph{Algebra and Logic} \textbf{28}, no.~3 (1989), 220--237. The supplied record refers to that paper for definitions but does not reproduce them. Thus the displayed convergence condition is a reference to an unavailable definition, not a specified topology or convergence criterion. The graph conventions, any local finiteness requirement, and any rooting or other auxiliary structure required by that definition remain unresolved. The interpretation of primitivity as the vertex action also requires confirmation from that reference.
\subsection{Short English statement}
Characterize the infinite connected graphs that arise as limits of sequences of finite connected graphs with primitive automorphism groups, using Trofimov's cited definitions. The supplied record omits those definitions, so the convergence notion and associated graph conventions require clarification.
\subsection{Sources}
\begin{itemize}
\item Local UnsolvedMath registry, supplied complete source record: permanent ID \(700906\), exact source problem code \texttt{catalogue entry}, title \emph{Kourovka Notebook Problem 12.89}. The registry labels the problem open; this is supplied source metadata, not an independently verified assessment.
\item \emph{The Kourovka Notebook (Unsolved Problems in Group Theory)}, original problem 12.89, 12th issue (1992), proposed by V. I. Trofimov. The supplied background also identifies the 21st issue edition (2026), \texttt{arXiv:1401.0300v48}. The collection metadata names E. I. Khukhro and V. D. Mazurov as editors.
\item V. I. Trofimov, \emph{Algebra and Logic} \textbf{28}, no.~3 (1989), 220--237. The problem explicitly directs the reader to this reference for definitions. No article title or source URL is supplied, and its contents have not been independently checked here.
\end{itemize}
\end{document}
